%------------------------------------------------------------------------------%
\section{Cálculo de Funções Reais}
\label{sec:Derivacao_Integracao}

Apresentamos nessa seção algumas definições e propriedades do Cálculo de
funções reais que podem ser úteis para o estudo dos conteúdos desse material.
Começamos pela definição de uma função real.

%------------------------------------------------------------------------------%
\begin{definicao}{Função}{def:funcao}
  Uma \emph{Função}\index{Função} $f$ de um conjunto $D$ para um conjunto $E$
  é uma regra que associa, sem ambiguidade, um único elemento $y=f(x)\in E$
  para cada elemento $x\in D$. Denotamos uma função por
  \[
    f: D \to E
  \]
\end{definicao}
%------------------------------------------------------------------------------%

Dizemos que $D$ é o \emph{Domínio}\index{Domínio} e $E$ o
\emph{Contra Domínio}\index{Contra Domínio} de $f$, enquanto o
conjunto de todos os valores $y=f(x)$ correspondentes a algum $x\in D$ é
a \emph{Imagem}\index{Imagem} de $f$.

%------------------------------------------------------------------------------%
\begin{definicao}{Função Real}{def:funcao-real}
  Uma \emph{Função Real}\index{Função!real} é uma função cujo domínio e
  contra domínio são subconjuntos dos reais $\R$.
\end{definicao}
%------------------------------------------------------------------------------%

A seguir apresentamos a definição formal do limite de uma função real.

%------------------------------------------------------------------------------%
\begin{definicao}{Limite de Função Real}{def:limite-funcao-real}
  Seja $f$ uma função definida em intervalo aberto que contém do ponto $a\in\R$,
  podendo não ser definida em $a$. Dizemos que o \emph{Limite}\index{Limite!função real}
  de $f(x)$ quando $x$ se aproxima de $a$ é $L\in\R$ e escrevemos
  \[
    \lim_{x\to a} f(x) = L
  \]
  se para cada $\varepsilon>0$ existir $\delta>0$ tal que
  \[
    0 < \abs{x-a} < \delta \qquad\Rightarrow\qquad \abs{f(x)-L} < \varepsilon
  \]
\end{definicao}
%------------------------------------------------------------------------------%

Essa definição pode ser estendida para o caso em que $a$ é $\pm\infty$.
Quando o limite existe ele possui as propriedades descritas na proposição a seguir.

%------------------------------------------------------------------------------%
\begin{proposicao}{Propriedades do Limite de Funções Reais}{pro:limite-funcao-real}
  Supondo que $c$ é uma constante, $f$ e $g$ são funções reais e que os limites
  \[
    \lim_{x\to a} f(x) \qquad\text{ e }\qquad \lim_{x\to a} g(x)
  \]
  existem, então valem as \emph{Propriedades do Limite}
  \begin{enumerate}
    \item \( \lim_{x\to a} \Big[ f(x) + g(x) \Big] = \lim_{x\to a} f(x) + \lim_{x\to a} g(x) \)
    \item \( \lim_{x\to a} \Big[ f(x) - g(x) \Big] = \lim_{x\to a} f(x) - \lim_{x\to a} g(x) \)
    \item \( \lim_{x\to a} cf(x) = c \lim_{x\to a} f(x) \)
    \item \( \lim_{x\to a} \Big[ f(x)g(x) \Big] = \lim_{x\to a} f(x) \cdot \lim_{x\to a} g(x) \)
    \item \( \lim_{x\to a} \frac{f(x)}{g(x)} = \dfrac{\displaystyle\lim_{x\to a} f(x)}
                                                     {\displaystyle\lim_{x\to a} d(x)} \)
          \qquad desde que \(\lim_{x\to a} d(x) \neq 0\)
  \end{enumerate}
\end{proposicao}
%------------------------------------------------------------------------------%

Uma propriedade de algumas funções reais é não possuir buracos ou saltos em seu
gráfico. Dizemos que essa funções são contínuas e a próxima definição nos fornece
uma caracterização para essas funções.

%------------------------------------------------------------------------------%
\begin{definicao}{Continuidade de Função Real}{def:continuidade-funcao-real}
  Uma função real é \emph{Contínua}\index{Função!contínua} em um ponto $a$ se
  \begin{enumerate}
    \item \( f(a) \) existe
    \item \( \lim_{x\to a} f(x) \) existe
    \item \( \lim_{x\to a} f(x) = f(a) \)
  \end{enumerate}
\end{definicao}
%------------------------------------------------------------------------------%

Um resultado importante sobre funções contínuas é o Teorema do Valor Intermediário.

%------------------------------------------------------------------------------%
\begin{teorema}{Teorema do Valor Intermediário}{teo:valor-intermediario}
  Suponha\index{Teorema!valor intermediário}
  que a função $f$ é contínua em todos os pontos de um intervalo fechado
  $[a, b]$ e que \(f(a)\neq f(b)\). Assuma que $p$ é um valor entre $f(a)$ e $f(b)$
  então existe um número $c\in(a, b)$ tal que $f(c)=p$.
\end{teorema}
%------------------------------------------------------------------------------%

Com o limite podemos definir a derivada de uma função real que nos fornece o
valor da inclinação da reta tangente ao gráfico da função em cada ponto.

%------------------------------------------------------------------------------%
\begin{definicao}{Derivada de Função Real}{def:derivada-funcao-real}
  A \emph{Derivada}\index{Derivada} de uma função real $f(x)$ em relação
  a variável $x$ é a função $f'$ cujo valor em $x$ é
  \[
    f'(x) = \lim_{h\to 0} \frac{\,f(x+h)-f(x)\,}{h}
  \]
  desde que o limite exista.
\end{definicao}
%------------------------------------------------------------------------------%

Ao calcularmos derivadas usamos as regras de derivação listadas na proposição
a seguir.

%------------------------------------------------------------------------------%
\begin{proposicao}{Regras de Derivação}{pro:regras-derivacao}
  Sejam $f$ e $g$ funções deriváveis, $c$ e $r$ constantes, então temos que
  \begin{multicols}{2}
    \begin{enumerate}
      \item \( \frac{d}{dx} c = 0 \)
      \item \( \frac{d}{dx} x^n = nx^{n-1} \)
      \item \( \frac{d}{dx} e^x = e^x \)
      \item \( \big(cf\big)' = cf' \)
      \item \( \big(f+g\big)' = f' + g' \)
      \item \( \big(f-g\big)' = f' - g' \)
      \item \( \big(fg\big)' = f'g + fg' \)
      \item \( \left(\frac{f}{\,g\,}\right)' = \frac{f'g - fg'}{g^2} \)
      \item Se \(h(x) = f\big(g(x)\big) \) então  \[ h'(x) = f'\big(g(x)\big)\,g'(x) \]
    \end{enumerate}
  \end{multicols}
\end{proposicao}
%------------------------------------------------------------------------------%

Uma propriedade importante de funções deriváveis é o Teorema de Rolle ou
Teorema do Valor Médio.

%------------------------------------------------------------------------------%
\begin{teorema}{Teorema de Rolle}{teo:rolle}
  Suponha\index{Teorema!de Rolle}\index{Teorema!valor médio}
  que $f$ é uma função
  \begin{enumerate}
    \item contínua  no intervalo fechado $[a, b]$
    \item derivável no intervalo aberto  $(a, b)$ e
    \item \(f(a)=f(b)\)
  \end{enumerate}
  então existe $c\in(a, b)$ onde $f'(c)=0$.
\end{teorema}
%------------------------------------------------------------------------------%

A operação inversa da derivação e a integração e é apresentadas nos próximos resultados.

%------------------------------------------------------------------------------%
\begin{definicao}{Primitiva de Função Real}{def:primitiva-funcao-real}
  A \emph{Primitiva}\index{Primitiva},
  \emph{Antiderivada}\index{Antiderivada} ou
  \emph{Integral Indefinida}\index{Integral!indefinida}
  de uma função $f$ é a função $F$ tal que
  \[
    \frac{dF}{dx}(x) = f(x)
  \]
  Denotamos a primitiva por
  \[
    F(x) = \int f(x)\, dx
  \]
\end{definicao}
%------------------------------------------------------------------------------%

É importante observar que a primitiva não é única, se $F(x)$ é uma primitiva de $f(x)$
então $F(x)+c$ também será para qualquer constante $c$. A proposição a seguir
apresenta algumas propriedades da primitiva.

%------------------------------------------------------------------------------%
\begin{proposicao}{Propriedades da Primitiva}{pro:propriedades-primitiva}
  Dadas funções integráveis $f$ e $g$ e a constante $c$ temos que
  \begin{enumerate}
    \item \( \int f(x)+g(x)\,dx = \int f(x)\,dx + \int g(x)\,dx\)
    \item \( \int f(x)-g(x)\,dx = \int f(x)\,dx - \int g(x)\,dx\)
    \item \( \int cf(x)\,dx   = c \int f(x)\,dx \)
  \end{enumerate}
\end{proposicao}
%------------------------------------------------------------------------------%

Uma operação relacionada é o cálculo da área entre o gráfico da função $f$ e
o eixo-$x$, chamamos essa operação de integral definida como descrito na
próxima definição.

%------------------------------------------------------------------------------%
\begin{definicao}{Integral Definida}{def:integral-defnida}
  A \emph{Integral Definida}\index{Integral!definida} de uma função real $f$
  é a área, considerando o sinal, entre o gráfico da função e o eixo-$x$
  no intervalo $[a, b]$
  \[
    \int_{a}^{b} f(x) \, dx
  \]
  Essa área é o limite da \emph{Soma de Riemman} correspondente.
\end{definicao}
%------------------------------------------------------------------------------%

A integral definida possui as propriedades descrita na proposição a seguir.

%------------------------------------------------------------------------------%
\begin{proposicao}{Propriedades da Integral Definida}{pro:propriedade-integral-definida}
  Dadas funções integráveis $f$ e $g$ e a constante $c$ temos que
  \begin{enumerate}
    \item \( \int_a^b c\,dx = c(b-a) \)
    \item \( \int_a^b f(x)+g(x)\,dx = \int_a^b f(x)\,dx + \int_a^b g(x)\,dx\)
    \item \( \int_a^b f(x)-g(x)\,dx = \int_a^b f(x)\,dx - \int_a^b g(x)\,dx\)
    \item \( \int_a^b cf(x)\,dx = c \int_a^b f(x)\,dx \)
    \item \( \int_a^c f(x)\,dx + \int_c^b f(x)\,dx = \int_a^b f(x)\,dx \)
  \end{enumerate}
\end{proposicao}
%------------------------------------------------------------------------------%

%------------------------------------------------------------------------------%
\begin{proposicao}{Integração por Partes}{pro:integracao-partes}
  Dadas funções integráveis $p$ e $q$, podemos escrever a seguinte relação
  entre as integrais indefinidas
  \[
    \int u\,dv = uv - \int v\,du \,.
  \]
  De modo equivalente temos a relação entre as integrais definidas
  \[
    \int_{-\infty}^\infty p\,dq = pq \, \evalat{-\infty}{\infty} - \int_{-\infty}^\infty q\,dp \,.
  \]
\end{proposicao}
%------------------------------------------------------------------------------%

O Teorema Fundamental do Cálculo é o resultado que determina as relações
entre a derivada e as integrais definida e indefinida de uma função.

%------------------------------------------------------------------------------%
\begin{teorema}{Teorema Fundamental do Cálculo}{teo:fundamental-calculo}
\begin{description}
  \item[Parte 1]\index{Teorema!fundamental do Cálculo}
  Se $f$ é uma função contínua em $[a, b]$ então a função $F$ definida por
  \[
    F(x) = \int_a^x f(t)\, dt \qquad\qquad a \leq x \leq b
  \]
  é contínua em $[a, b]$, diferenciável em $(a, b)$ e
  \[
    F'(x) = f(x)
  \]
  \item[Parte 2]
  Se $f$ é contínua em $[a, b]$ então
  \[
   \int_a^b f(x)\,dx = F(b) - F(a) = F(x)\evalat{a}{b}
  \]
  onde $F$ é qualquer primitiva de $f$.
\end{description}
\end{teorema}
%------------------------------------------------------------------------------%

As integrais definidas são calculadas em um intervalo finito e a função $f$
não deve ter nenhuma descontinuidade infinita. Porém, mesmo se o intervalo for infinito
ou $f$ possuir uma descontinuidade infinita, em alguns casos, é possível
calcular a área sob o gráfico usando a integral imprópria.

%------------------------------------------------------------------------------%
\begin{definicao}{Integral Imprópria -- Tipo 1}{def:integral-impropria-1}
  A \emph{Integral Imprópria}\index{Integral!imprópria} do tipo 1
  é a integral definida de uma função $f$ em um intervalo infinito.

  Se \( \int_a^b f(x)\, dx \)
  existe para todo $b\geq a$ e o limite existir, temos que
  \[
    \int_a^\infty f(x)\,dx = \lim_{b\to\infty}\int_a^b f(x)\, dx
  \]
  Se \( \int_a^b f(x)\, dx \)
  existe para todo $a\leq b$ e o limite existir, temos que
  \[
    \int_{-\infty}^b f(x)\,dx = \lim_{a\to-\infty}\int_a^b f(x)\, dx
  \]
  Se \( \int_a^\infty    f(x)\,dx \)
  e  \( \int_{-\infty}^b f(x)\,dx \)
  são convergentes, então
  \[
    \int_{-\infty}^\infty f(x)\,dx = \int_{-\infty}^c f(x)\, dx
                                   + \int_c^\infty    f(x)\, dx
  \]
\end{definicao}
%------------------------------------------------------------------------------%

%------------------------------------------------------------------------------%
\begin{definicao}{Integral Imprópria -- Tipo 2}{def:integral-impropria-2}
  Se $f$ é contínua em $[a, b)$ e descontínua em $b$, e o limite existir,
  temos que
  \[
    \int_a^b f(x)\,dx = \lim_{t\to b} \int_a^t f(x)\,dx
  \]
  Se $f$ é contínua em $(a, b]$ e descontínua em $a$, e o limite existir,
  temos que
  \[
    \int_a^b f(x)\,dx = \lim_{t\to a} \int_t^b f(x)\,dx
  \]
\end{definicao}
%------------------------------------------------------------------------------%

\input{A-Complementar/08-1-Derivadas}
%------------------------------------------------------------------------------%

\begin{nlongtable}[\setlength{\extrarowheight}{5mm}]%
                  {>{$\displaystyle\vphantom{\int_\int}\qquad}w{l}{0.300\linewidth}<{$}
                   >{$\displaystyle\vphantom{\int_\int}}      w{l}{0.682\linewidth}<{$}}
                  {Tabela de Integrais}
                  {tab:tabela-integrais}
  \int u \,dv                          & = uv - \int v \,du                                        \\ \hline
  \int u^n \,du                        & = \frac{u^{n+1}}{\,n+1\,} + C \qquad n\neq -1             \\ \hline
  \int \frac{du}{u}                    & = \ln\abs{u} + C                                          \\ \hline
  \int e^u \,du                        & = e^u + C                                                 \\ \hline
  \int a^u \,du                        & = \frac{a^u}{\,\ln(a)\,} + C                              \\ \hline
  \int \ln(u)\,du                      & = u\big( \ln(u) - 1 \big) + C                             \\ \hline
  \int \sin(u)   \,du                  & = -\cos(u) + C                                            \\ \hline
  \int \cos(u)   \,du                  & =  \sin(u) + C                                            \\ \hline
  \int \sec^2(u) \,du                  & =  \tan(u) + C                                            \\ \hline
  \int \csc^2(u) \,du                  & = -\cot(u)+ C                                             \\ \hline
  \int \sec(u) \tan(u) \,du            & =  \sec(u) + C                                            \\ \hline
  \int \csc(u) \cot(u) \,du            & = -\csc(u) + C                                            \\ \hline
  \int \tan(u) \,du                    & = \ln\abs{\sec(u)} + C                                    \\ \hline
  \int \cot(u) \,du                    & = \ln\abs{\sin(u)} + C                                    \\ \hline
  \int \sec(u) \,du                    & = \ln\abs{\sec(u)+\tan(u)} + C                            \\ \hline
  \int \csc(u) \,du                    & = \ln\abs{\csc(u)-\cot(u)} + C                            \\ \hline
  \int \frac{du}{\sqrt{a^2-u^2\,}\,}   & = \arcsin\left(\frac{\,u\,}{a}\right) + C                 \\ \hline
  \int \frac{du}{\,a^2+u^2\,}          & = \frac{1}{\,a\,} \arctan\left(\frac{\,u\,}{a}\right)+ C  \\ \hline
  \int \frac{du}{u\sqrt{u^2-a^2\,}\,}  & = \frac{1}{\,a\,} \arcsec\left(\frac{\,u\,}{a}\right) + C \\ \hline
  \int \frac{du}{\,a^2-u^2\,}          & = \frac{1}{\,2a\,}\ln\abs{\frac{\,u+a\,}{u-a}} + C        \\ \hline
  \int \frac{du}{\,u^2-a^2\,}          & = \frac{1}{\,2a\,}\ln\abs{\frac{\,u-a\,}{u+a}} + C        \\ \hline
  \int u \sin(u) \, du                 & = \sin(u) - u \cos(u) + C                                 \\ \hline
  \int u \cos(u) \, du                 & = \cos(u) + u \sin(u) + C                                 \\ \hline
  \int u^2 \sin(u) \, du               & = \hpm 2\cos(u) + 2u \sin(u) -u^2\cos(u) + C              \\ \hline
  \int u^2 \cos(u) \, du               & =    - 2\sin(u) + 2u \cos(u) +u^2\sin(u) + C              \\ \hline
  \int u^n \sin(u) \, du               & =    - u^n\cos(u) + n\int u^{n-1} \cos(u)\,du             \\ \hline
  \int u^n \cos(u) \, du               & = \hpm u^n\sin(u) - n\int u^{n-1} \sin(u)\,du
\end{nlongtable}

%------------------------------------------------------------------------------%


%------------------------------------------------------------------------------%
